\documentclass[10pt,a4paper]{article}
\usepackage{fontspec}
\setmainfont{Noto Sans CJK SC}
\setsansfont{Noto Sans CJK SC}
\setmonofont{DejaVu Sans Mono}
\XeTeXlinebreaklocale "zh"
\XeTeXlinebreakskip=0pt plus 1pt
\usepackage[margin=0.64in,top=0.62in,bottom=0.62in,headheight=15pt]{geometry}
\usepackage{graphicx,array,tabularx,colortbl,xcolor,fancyhdr,hyperref,fvextra}
\hypersetup{colorlinks=true,urlcolor=blue!40!black,linkcolor=black,pdftitle={嵌入式系统作业2},pdfauthor={}}
\definecolor{tablehead}{HTML}{DFE9F0}
\definecolor{tabline}{HTML}{D9D9D9}
\pagestyle{fancy}\fancyhf{}\fancyhead[L]{\small 嵌入式系统\quad |\quad 作业2}\fancyfoot[C]{\small\thepage}
\renewcommand{\headrulewidth}{0pt}
\setlength{\parindent}{0pt}\setlength{\parskip}{5pt}
\setlength{\emergencystretch}{2em}
\hyphenpenalty=10000\exhyphenpenalty=10000
\renewcommand{\baselinestretch}{1.10}
\setlength{\tabcolsep}{5pt}\renewcommand{\arraystretch}{1.20}
\arrayrulecolor{tabline}
\newcolumntype{Y}{>{\raggedright\arraybackslash}X}
\newcommand{\qtitle}[1]{{\Large\bfseries #1}\par\vspace{4pt}}
\newcommand{\qsub}[1]{{\bfseries #1}\par}
\begin{document}
\pdfbookmark[0]{嵌入式系统作业2}{page1}
\qtitle{嵌入式系统作业2}
{\small Cortex-M4体系结构与STM32基础  彭俊杰  智装备241  2024306100324}\par
本作业按课程题目依次说明Cortex-M4的存储器与异常机制、STM32F407的总线、时钟及复位电路，再结合教师CH5EX01例程与四份参考工程分析启动、GPIO和SysTick。器件参数分别按STM32F407ZGT6和STM32L496VET6核对；代码计算、历史编译记录和实际运行结果分别说明。

\qsub{第1题 英文缩写}
\begingroup\small
\begin{tabularx}{\linewidth}{|>{\hsize=0.335\hsize}Y|>{\hsize=1.0\hsize}Y|>{\hsize=1.665\hsize}Y|}\hline
\rowcolor{tablehead}
\textbf{缩写} & \textbf{英文名称} & \textbf{含义与作用} \\ \hline
ARM & Acorn RISC Machine；后称 Advanced RISC Machines & 历史名称；本课程指 Arm 处理器架构及其内核，Cortex-M4 是其中一种内核。 \\ \hline
AHB & Advanced High-performance Bus & 高级高性能总线，连接处理器、存储器和较高速外设。 \\ \hline
AHB-Lite & AHB Lite & AHB 的单主设备接口子集；系统仍可通过总线矩阵组合多个主设备。 \\ \hline
APB & Advanced Peripheral Bus & 高级外设总线，接口较简单，常连接控制类和较低速外设。 \\ \hline
PPB & Private Peripheral Bus & 私有外设总线；对应内核私有外设与调试部件的地址区域。 \\ \hline
DAP & Debug Access Port & 调试访问端口；允许调试器访问处理器与系统资源。 \\ \hline
JTAG & Joint Test Action Group & 联合测试工作组名称，通常指 IEEE 1149.1 测试与调试接口。 \\ \hline
SWD & Serial Wire Debug & 串行线调试，主要使用 SWDIO 与 SWCLK 两根信号线。 \\ \hline
MPU & Memory Protection Unit & 存储器保护单元；给地址区域设置访问权限与属性，不等同于地址转换MMU。 \\ \hline
FPU & Floating-Point Unit & 浮点运算单元；两款所选器件的 Cortex-M4F 支持单精度浮点。 \\ \hline
\end{tabularx}
\endgroup\par
术语依据：内核编程手册及厂商器件资料 [1–4]。DAP 是调试系统的组成部分，JTAG、SWD 是可采用的调试访问接口。


\clearpage
\pdfbookmark[0]{第1题 英文缩写续}{page2}
\qtitle{第1题 英文缩写续}
\begingroup\small
\begin{tabularx}{\linewidth}{|>{\hsize=0.335\hsize}Y|>{\hsize=1.0\hsize}Y|>{\hsize=1.665\hsize}Y|}\hline
\rowcolor{tablehead}
\textbf{缩写} & \textbf{英文名称} & \textbf{含义与作用} \\ \hline
NVIC & Nested Vectored Interrupt Controller & 嵌套向量中断控制器，管理外部中断的使能、挂起、优先级和抢占。 \\ \hline
ISR & Interrupt Service Routine & 中断服务程序，处理中断事件后返回原程序。 \\ \hline
Bit-band & Bit-banding & 位带：把支持区域中的一个位映射到别名区中的一个字地址，便于按位访问。 \\ \hline
SCS & System Control Space & 系统控制空间，包含系统控制块、NVIC、SysTick等寄存器。 \\ \hline
SysTick & System Tick Timer & 系统滴答定时器，24位递减计数器，可作为周期时基。 \\ \hline
MSP & Main Stack Pointer & 主栈指针，复位后使用，异常处理时使用MSP。 \\ \hline
PC & Program Counter & 程序计数器，R15；与当前执行和下一步取指位置有关。 \\ \hline
SP & Stack Pointer & 栈指针，R13；可指向MSP或PSP所管理的栈。 \\ \hline
GPIO & General-Purpose Input/\allowbreak{}Output & 通用输入输出端口，可配置输入、输出、复用和模拟模式。 \\ \hline
TIM & Timer & 定时器；可用于定时、计数、捕获、比较和PWM输出。 \\ \hline
UART & Universal Asynchronous Receiver/\allowbreak{}Transmitter & 通用异步收发器；通过TX、RX等引脚收发串行数据。 \\ \hline
CAN Bus & Controller Area Network Bus & 控制器局域网络总线；常用于节点间可靠通信，芯片CAN控制器通常还需外部收发器。 \\ \hline
\end{tabularx}
\endgroup\par
\qsub{几组容易混淆的概念}
NVIC是硬件控制器，ISR是软件处理程序；SysTick是内核定时器，TIM2是芯片外设定时器。UART只强调异步通信，USART还可支持同步方式，因此 USART1 的名称不能直接改写成 UART1。

地址区域由体系结构规定，区域中具体实现哪些存储器和外设由芯片决定。不能看到“SRAM区有512 MB地址空间”，就认为芯片装有512 MB SRAM。


\clearpage
\pdfbookmark[0]{第2题 存储器映射与外设地址}{page3}
\qtitle{第2题 存储器映射与外设地址}
Cortex-M4 使用32位地址，统一编址范围为 0x00000000 至 0xFFFFFFFF，共4 GB。指令、数据和寄存器都可通过地址定位，主要分区如下 [2–3]。

\begin{center}
\includegraphics[width=6.3in]{assets/图1_CortexM4存储器映射.png}\par
{\footnotesize 图1 Cortex-M4存储器映射示意图 地址框大小不按容量比例绘制}
\end{center}
\begingroup\small
\begin{tabularx}{\linewidth}{|>{\hsize=0.674\hsize}Y|>{\hsize=1.186\hsize}Y|>{\hsize=1.14\hsize}Y|}\hline
\rowcolor{tablehead}
\textbf{对象} & \textbf{地址或范围} & \textbf{所在区域} \\ \hline
GPIOA /\allowbreak{} GPIOF & 0x40020000 /\allowbreak{} 0x40021400 & 片上外设区，AHB1 \\ \hline
CAN1 /\allowbreak{} CAN2 & 0x40006400 /\allowbreak{} 0x40006800 & 片上外设区，APB1 \\ \hline
SysTick & 0xE000E010—0xE000E01C & 系统区中的SCS，属于内核私有外设 \\ \hline
F407主Flash & 0x08000000—0x080FFFFF & 代码区，ZGT6容量1 MB \\ \hline
F407 SRAM1和SRAM2 & 0x20000000—0x2001FFFF & SRAM区，合计128 KB \\ \hline
F407 CCM数据RAM & 0x10000000—0x1000FFFF & 落在代码区地址范围，容量64 KB \\ \hline
\end{tabularx}
\endgroup\par
位带补充：SRAM位带区前1 MB位于0x20000000起，别名区从0x22000000起；外设位带区从0x40000000起，别名区从0x42000000起。位带功能和具体可访问地址均须按所选芯片核对。CCM并不因为名字含RAM就位于0x20000000区，也不能直接用作F407通用DMA缓冲区。


\clearpage
\pdfbookmark[0]{第3题 中断与异常向量表}{page4}
\qtitle{第3题 中断与异常向量表}
中断使处理器在事件发生时暂时转去执行处理程序，处理后再返回。Cortex-M4把复位、故障、系统服务和外部中断统一纳入异常机制；外部中断只是异常的一类 [3]。

\begin{center}
\includegraphics[width=6.3in]{assets/图2_中断响应过程.png}\par
{\footnotesize 图2 外部中断从请求到返回的处理过程}
\end{center}
关键词：中断源、请求、使能、屏蔽、挂起、激活、优先级、抢占、嵌套、现场保护、ISR、清除标志、异常返回。挂起表示请求已记录，激活表示正在服务；同一中断也可同时处于激活和挂起状态。可配置优先级通常是数值越小、优先级越高。

\qsub{向量表为什么能找到处理程序}
向量表是一张入口地址表，每项占4字节。第0项保存初始MSP，后续项保存各异常处理函数地址。硬件根据异常编号定位入口，不需要程序逐个判断所有中断源。复位向量不是初始栈顶。

\begingroup\small
\begin{tabularx}{\linewidth}{|>{\hsize=0.558\hsize}Y|>{\hsize=0.93\hsize}Y|>{\hsize=1.512\hsize}Y|}\hline
\rowcolor{tablehead}
\textbf{向量表偏移} & \textbf{项目} & \textbf{作用} \\ \hline
0x00 & 初始MSP & 建立主栈的起始栈顶 \\ \hline
0x04 & Reset\_\allowbreak{}Handler & 复位后执行的入口 \\ \hline
0x08 /\allowbreak{} 0x0C & NMI /\allowbreak{} HardFault & 不可屏蔽中断 /\allowbreak{} 硬故障 \\ \hline
0x3C & SysTick\_\allowbreak{}Handler & 异常号15的系统滴答入口 \\ \hline
0x40起 & 外部IRQ向量 & 异常号为IRQn+16 \\ \hline
0xB0 & TIM2\_\allowbreak{}IRQHandler & TIM2\_\allowbreak{}IRQn=28，偏移4×(16+28) \\ \hline
\end{tabularx}
\endgroup\par
异常进入时，硬件基本栈帧保存R0—R3、R12、LR、PC和xPSR；浮点上下文可能增加额外保存。ISR应按外设规则清除中断源，否则返回后可能再次进入。VTOR可设置向量表基址，搬移时必须满足对齐与地址要求。


\clearpage
\pdfbookmark[0]{第4题 STM32F407片上总线}{page5}
\qtitle{第4题 STM32F407片上总线}
总线矩阵的主、从端口与AHB1、AHB2、AHB3、APB1、APB2这组外设总线名称属于不同层次。STM32F407采用多层AHB矩阵连接8个主端口与7个从端口；下页列出矩阵端口。本页先回答GPIO、TIM2与USART1挂在哪条外设总线上 [2，RM0090第2.1节]。

\begin{center}
\includegraphics[width=6.3in]{assets/图3_F407片上总线.png}\par
{\footnotesize 图3 STM32F407片上外设总线的简化连接}
\end{center}
\begingroup\small
\begin{tabularx}{\linewidth}{|>{\hsize=0.581\hsize}Y|>{\hsize=0.465\hsize}Y|>{\hsize=1.953\hsize}Y|}\hline
\rowcolor{tablehead}
\textbf{模块} & \textbf{所属总线} & \textbf{使能位置与基址} \\ \hline
GPIOA—GPIOI & AHB1 & RCC\_\allowbreak{}AHB1ENR；GPIOA基址0x40020000 \\ \hline
TIM2 & APB1 & RCC\_\allowbreak{}APB1ENR.TIM2EN；基址0x40000000 \\ \hline
USART1 & APB2 & RCC\_\allowbreak{}APB2ENR.USART1EN；基址0x40011000 \\ \hline
\end{tabularx}
\endgroup\par
\qsub{从连接关系理解初始化}
访问外设前先开启其所属总线上的外设时钟。例如配置GPIOF，需要使能AHB1上的GPIOF时钟；配置TIM2和USART1则分别使能APB1和APB2中的对应位。芯片已运行不代表每个外设时钟都已打开。

AHB1上的GPIO寄存器通过总线访问，但引脚电平变化还受GPIO输出速度配置、负载、电气条件和软件访问方式影响，不能把168 MHz总线频率直接写成GPIO输出方波频率。

AHB-Lite是单主接口协议，不意味着整个芯片只能有一个主设备。多层矩阵把多个这样的主接口连接到不同从端口，允许对不同目标的访问并行进行。


\clearpage
\pdfbookmark[0]{第4题 总线矩阵的主从端口}{page6}
\qtitle{第4题 总线矩阵的主从端口}
总线矩阵的主端口发起访问，从端口响应访问。以下名称按RM0090 Rev19第2.1节及其系统结构图列出，避免把一个DMA控制器只算作一个端口，或把AHB到APB桥重复算成独立矩阵从端口。

\begingroup\small
\begin{tabularx}{\linewidth}{|>{\hsize=0.769\hsize}Y|>{\hsize=1.231\hsize}Y|}\hline
\rowcolor{tablehead}
\textbf{主端口} & \textbf{访问用途} \\ \hline
1 Cortex-M4 I-bus & 指令取指 \\ \hline
2 Cortex-M4 D-bus & 代码区的数据访问 \\ \hline
3 Cortex-M4 S-bus & 系统存储器与外设访问 \\ \hline
4 DMA1 memory bus & DMA1存储器访问 \\ \hline
5 DMA2 memory bus & DMA2存储器访问 \\ \hline
6 DMA2 peripheral bus & DMA2外设访问 \\ \hline
7 Ethernet DMA & 以太网DMA访问 \\ \hline
8 USB OTG HS DMA & USB高速控制器DMA访问 \\ \hline
\end{tabularx}
\endgroup\par
\begingroup\small
\begin{tabularx}{\linewidth}{|>{\hsize=0.769\hsize}Y|>{\hsize=1.231\hsize}Y|}\hline
\rowcolor{tablehead}
\textbf{从端口} & \textbf{响应目标} \\ \hline
1 Flash ICode & Flash指令接口 \\ \hline
2 Flash DCode & Flash数据接口 \\ \hline
3 SRAM1 & 112 KB SRAM \\ \hline
4 SRAM2 & 16 KB SRAM \\ \hline
5 AHB1 peripherals & AHB1外设，包含AHB到APB桥及其后方外设 \\ \hline
6 AHB2 peripherals & AHB2外设 \\ \hline
7 FSMC & 外部存储器控制器 \\ \hline
\end{tabularx}
\endgroup\par
64 KB CCM数据RAM不在上述矩阵中，由内核的专用路径访问，因此不能把CCM加入这7个从端口，也不能把它当作普通DMA可达的SRAM。GPIO仍属于AHB1，TIM2属于APB1，USART1属于APB2。


\clearpage
\pdfbookmark[0]{第5题 时钟源与系统时钟}{page7}
\qtitle{第5题 时钟源与系统时钟}
\begingroup\small
\begin{tabularx}{\linewidth}{|>{\hsize=0.465\hsize}Y|>{\hsize=0.302\hsize}Y|>{\hsize=2.233\hsize}Y|}\hline
\rowcolor{tablehead}
\textbf{类别} & \textbf{时钟源} & \textbf{标称值或范围与用途} \\ \hline
高速内部 & HSI & 16 MHz RC，可直接作为SYSCLK或作为主PLL输入 \\ \hline
高速外部 & HSE & 晶体/\allowbreak{}陶瓷谐振器通常4—26 MHz；外部有源时钟走旁路方式并核对电气限制 \\ \hline
低速内部 & LSI & 约32 kHz RC，主要供IWDG，也可供RTC \\ \hline
低速外部 & LSE & 32.768 kHz晶体，常用于RTC \\ \hline
\end{tabularx}
\endgroup\par
\begin{center}
\includegraphics[width=6.3in]{assets/图4_F407时钟路径.png}\par
{\footnotesize 图4 STM32F407的主要时钟来源与分配路径}
\end{center}
SYSCLK可选择HSI、HSE或主PLL输出。主PLL以HSI/\allowbreak{}HSE为输入，先经M分频、N倍频，再经P分频，得到系统候选时钟：fPLL＝f输入÷M×N÷P。RCC\_\allowbreak{}CFGR中的SW发出切换请求，SWS反映实际已切换到的来源 [2]。

上限：AHB1所属HCLK最高168 MHz，APB1的PCLK1最高42 MHz，APB2的PCLK2最高84 MHz。达到最高频率还要满足供电、稳压等级和Flash等待周期等条件，不能只修改PLL参数。

计算例：假设HSE＝8 MHz，取M＝8、N＝336、P＝2，则VCO输入1 MHz、VCO输出336 MHz、SYSCLK＝168 MHz。AHB不分频，APB1四分频，APB2二分频，即HCLK＝168、PCLK1＝42、PCLK2＝84 MHz；Q＝7可得到48 MHz输出。

TIM2挂在APB1，但本例APB1分频不为1，故TIM2内核时钟为2×42＝84 MHz。上电复位后默认使用HSI，168 MHz并不是芯片无需配置就自动得到的工作频率。

电源条件须另行核对：STM32F407的VDD应用范围为1.8至3.6 V。开发板的USB/\allowbreak{}5 V入口经过稳压器后才供给芯片；不能把板级输入电压或某些5 V容忍I/\allowbreak{}O的能力当成VDD额定值。


\clearpage
\pdfbookmark[0]{第6题 复位引脚与典型电路}{page8}
\qtitle{第6题 复位引脚与典型电路}
STM32F407ZGT6中的Z表示144引脚，T表示LQFP封装。查DS8626 Rev12的LQFP144顶视图和引脚表可知，NRST在第25脚；第14脚是PF4。LQFP100封装的NRST才是第14脚，不能混用 [1，图14及表7]。

\begin{center}
\includegraphics[width=6.3in]{assets/图5_NRST复位电路.png}\par
{\footnotesize 图5 STM32F407ZGT6外部复位电路 依据厂家建议连接方式绘制}
\end{center}
\qsub{名称含义与动作过程}
NRST中的N表示低电平有效，RST表示Reset。平时上拉使NRST保持高电平，芯片正常运行；按下常开按键时NRST被拉至地，触发系统复位；松开后电容经上拉路径充电，NRST恢复高电平，处理器重新执行复位启动过程。

芯片内部已有弱上拉，外接10 kΩ上拉是常见设计选择，并非所有板卡都必须外加。100 nF电容应靠近NRST与地，帮助抑制干扰；具体复位脉宽和电平阈值按数据手册核验。图中的10 kΩ与100 nF名义乘积为1 ms，但这不等于准确的复位保持时间。

调试器的NRST可接入同一复位网络，驱动方式应避免与芯片或其他复位器件争用。除了外部引脚复位，芯片还有上电、掉电、欠压、看门狗和软件等复位来源；复位后可读取RCC中的复位原因标志帮助定位问题。

硬件检查顺序：先确认芯片完整丝印与封装，再按顶视图找到1脚方向，沿引脚编号定位25脚，最后用原理图网络名核对。不要按开发板排针编号代替芯片引脚编号。


\clearpage
\pdfbookmark[0]{第7题 两款主芯片的数据手册与特征}{page9}
\qtitle{第7题 两款主芯片的数据手册与特征}
本题采用STM32F407ZGT6和STM32L496VET6。L496VET6是本次课程题目指定型号；若实物板丝印不同，应另查对应器件手册，不能直接套用参数。两个官方PDF均随作业保留在sources目录。

\begingroup\small
\begin{tabularx}{\linewidth}{|>{\hsize=0.488\hsize}Y|>{\hsize=1.256\hsize}Y|>{\hsize=1.256\hsize}Y|}\hline
\rowcolor{tablehead}
\textbf{比较项目} & \textbf{STM32F407ZGT6} & \textbf{STM32L496VET6} \\ \hline
内核 & Cortex-M4，单精度FPU，DSP，MPU & Cortex-M4，单精度FPU，DSP，MPU \\ \hline
最高主频 & 168 MHz & 80 MHz \\ \hline
Flash & 1 MB，型号容量码G & 512 KB，型号容量码E \\ \hline
SRAM & 192 KB，含64 KB CCM；另有4 KB备份SRAM & 320 KB，含64 KB带硬件奇偶校验的SRAM2 \\ \hline
封装 & LQFP144，Z=144脚 & LQFP100，V=100脚 \\ \hline
温度等级6 & −40至+85 ℃ & −40至+85 ℃ \\ \hline
主电源范围 & VDD 1.8—3.6 V & 使用内部LDO时VDD 1.71—3.6 V \\ \hline
模拟外设 & 3个12位ADC，2路12位DAC & 3个12位ADC，2路12位DAC，运放、比较器 \\ \hline
通信与扩展 & CAN、USART/\allowbreak{}UART、SPI、I²C、USB OTG、以太网MAC、FSMC等 & CAN、USART/\allowbreak{}UART/\allowbreak{}LPUART、SPI、I²C、USB OTG FS、QUADSPI、LCD等 \\ \hline
主要特点 & 较高主频，网络接口与外部存储扩展丰富 & 更大SRAM容量，多种低功耗模式及低功耗外设 \\ \hline
\end{tabularx}
\endgroup\par
F407的以太网模块是MAC，实际连网还需要PHY与相关电路；其USB高速功能需要外部ULPI PHY。L496的数据手册标题写“最高1 MB Flash”是系列上限，而VET6的实际Flash为512 KB。外设是否引出到板上还要查封装复用表与板级原理图 [1，4]。

\href{https://www.st.com/resource/en/datasheet/stm32f407zg.pdf}{STM32F407ZG官方数据手册 DS8626}\par
\href{https://www.st.com/resource/en/datasheet/stm32l496ve.pdf}{STM32L496VE官方数据手册 DS11585}\par
阅读顺序：先读Features和器件对照表，核对完整型号；再读Pinouts和Memory map；最后结合具体外设查参考手册。数据手册主要给参数、引脚和电气条件，参考手册给寄存器与使用流程。


\clearpage
\pdfbookmark[0]{第8题 扩展练习一 最小系统}{page10}
\qtitle{第8题 扩展练习一 最小系统}
最小系统至少应具备稳定供电、去耦、复位、启动电平与下载调试通路。外部晶体并非启动程序的必要条件，可先用内部HSI检查供电、复位和SWD；需要准确频率时再配置HSE。

\begin{center}
\includegraphics[width=6.3in]{assets/图6_F407最小系统组成.png}\par
{\footnotesize 图6 STM32F407ZGT6最小系统组成 用于读图与核对连接}
\end{center}
VDD、VDDA、VSS、VSSA等同类引脚要逐一连接，不能只接其中一对。VCAP1和VCAP2是内部稳压器的电容端，各按手册配置2.2 μF低ESR电容，不作为外部负载电源。BOOT0下拉用于选择主Flash启动；SWD保留PA13、PA14和公共地，建议同时引出NRST与目标电压参考。

\qsub{课程原理图来源}
采用课程资料提供的普中F407-T100开发板原理图，主芯片为STM32F407ZGT6。图6是依据器件手册整理的最小系统连接示意；下一页给出原理图局部，assets目录保留完整图。两种图各自注明性质，便于从系统组成转到实际网络连接。

\href{https://drive.google.com/file/d/1t4c7yAtB6ed2Qf-TAu8f8MEismKaucLW/view}{课程资料 普中F407-T100开发板原理图}\par

\clearpage
\pdfbookmark[0]{第8题 扩展练习一 原理图局部}{page11}
\qtitle{第8题 扩展练习一 原理图局部}
下图摘自课程提供的普中F407-T100开发板原理图，保留原图标注。用它对照芯片引脚、电源和时钟相关网络；完整单页图另附在assets目录。来源按课程资料标注。

\begin{center}
\includegraphics[width=6.3in]{assets/课程资料_主芯片电路局部.png}\par
{\footnotesize 图7 课程资料中普中F407-T100主芯片电路局部}
\end{center}
读图时先按网络名追踪VDD、VSS、VDDA、VSSA、VCAP、BOOT0和NRST，再检查HSE与LSE连接。LED0与KEY0的网络关系用于后面的GPIO练习，不能由芯片型号推断所有开发板都采用相同接法。

\href{https://drive.google.com/file/d/1t4c7yAtB6ed2Qf-TAu8f8MEismKaucLW/view}{原始课程文件 普中F407-T100开发板原理图}\par

\clearpage
\pdfbookmark[0]{第8题 扩展练习二 实际工程结构}{page12}
\qtitle{第8题 扩展练习二 实际工程结构}
标准外设库SPL和硬件抽象库HAL都在寄存器之上提供接口。本节先说明已有四份参考工程，再单列教师CH5EX01 HAL/\allowbreak{}STD。它们使用同一类芯片，却不采用相同的板级引脚、晶振和延时实现，不能拼接成一个目标。

\begingroup\small
\begin{tabularx}{\linewidth}{|>{\hsize=0.558\hsize}Y|>{\hsize=1.233\hsize}Y|>{\hsize=1.209\hsize}Y|}\hline
\rowcolor{tablehead}
\textbf{工程} & \textbf{实际分组或关键文件} & \textbf{本次阅读重点} \\ \hline
SPL时钟系统 & User、Startup、StdPeriph\_\allowbreak{}Driver、CMSIS、APP & main中调用RCC\_\allowbreak{}HSE\_\allowbreak{}Config；PF9输出翻转 \\ \hline
SPL SysTick & 除上述分组外增加Public；SysTick.c位于Public & SystemInit设置时钟；轮询SysTick完成延时 \\ \hline
HAL最小工程 & User、CMSIS、STM32F4xx\_\allowbreak{}HAL\_\allowbreak{}Driver & main、HAL时基、PF9/\allowbreak{}PF10输出与PE4输入 \\ \hline
CubeMX生成工程 & Application/\allowbreak{}MDK-ARM、Application/\allowbreak{}User、Drivers及::CMSIS & clock168\_\allowbreak{}gpio.ioc、Src/\allowbreak{}main.c、Src/\allowbreak{}gpio.c \\ \hline
\end{tabularx}
\endgroup\par
\qsub{模板固定了什么}
SPL器件宏为STM32F40\_\allowbreak{}41xxx，库开关为USE\_\allowbreak{}STDPERIPH\_\allowbreak{}DRIVER。HAL最小工程的宏为USE\_\allowbreak{}HAL\_\allowbreak{}DRIVER、STM32F407xx和HSE\_\allowbreak{}VALUE=8000000。CubeMX目标使用USE\_\allowbreak{}HAL\_\allowbreak{}DRIVER、STM32F407xx，晶振设置需对照IOC及配置头文件。宏、启动文件、Include路径和驱动版本必须配套。

Keil中的group用于组织工程文件，不必与磁盘目录一一对应。例如HAL工程把启动文件放在CMSIS组中，不能仅凭组名判断它是否存在。真正参与构建的文件由.uvprojx目标配置决定。

\qsub{本次环境与结果范围}
上传的历史构建记录中的环境为µVision 5.35、Arm Compiler 5.06 update 7 build 960、STM32F4xx\_\allowbreak{}DFP 2.16.0；CubeMX为6.18.1，STM32CubeF4固件包为V1.28.3。构建结果依据2026年10月4日上传的原始日志，详细结果见后文。

没有实体开发板。现有证据包括源码、上传的AC5历史构建日志与仿真受阻截图；本轮未重新运行Keil目标编译，也未完成GPIO、PLL或SysTick的板上测量。修订源码的正确性还需在匹配工具链和硬件条件下验证。


\clearpage
\pdfbookmark[0]{第8题 扩展练习三 启动文件与复位序列}{page13}
\qtitle{第8题 扩展练习三 启动文件与复位序列}
实际HAL工程使用startup\_\allowbreak{}stm32f407xx.s。图8直接摘录该文件，保留源文件行号；省略号只表示未在图中展示的行，真实工程中的向量项没有删除。

\begin{center}
\includegraphics[width=6.3in]{assets/图8_真实启动文件.png}\par
{\footnotesize 图8 实际HAL启动源码摘录 第33 44 61至76和170至178行选录}
\end{center}
\qsub{从复位到main}
硬件先读取向量表第0项作为初始MSP，再读取第1项进入Reset\_\allowbreak{}Handler。图中DCD放置32位地址；第61、62行分别对应初始栈顶和复位入口。SysTick向量位于第76行，仍是入口地址而不是处理函数体。

第175至178行先调用SystemInit，再跳入Arm C运行库入口\_\allowbreak{}\_\allowbreak{}main，完成数据段初始化和C运行环境准备后进入main。该HAL工程的SystemInit主要配置FPU和VTOR；168 MHz的PLL配置在main调用的SystemClock\_\allowbreak{}Config中完成。两个普中参考SPL工程则在SystemInit内调用SetSysClock。

本文件Stack\_\allowbreak{}Size=0x400，即1 KB；Heap\_\allowbreak{}Size=0x200，即512 B。它们是本工程的栈、堆预留值，不是芯片固定容量。异常嵌套和较深的函数调用会增加栈需求，不能以“编译成功”代替栈余量评估。


\clearpage
\pdfbookmark[0]{第8题 扩展练习四 HAL GPIO初始化}{page14}
\qtitle{第8题 扩展练习四 HAL GPIO初始化}
HAL最小工程将PF9和PF10配置为推挽输出，将PE4配置为上拉输入。课程普中F407-T100接法中，LED低电平点亮、按键按下接地；换板时仍需核对原理图。下面摘录实际main.c中的初始化函数。

\begin{center}
\includegraphics[width=6.3in]{assets/图9_HAL_GPIO初始化.png}\par
{\footnotesize 图9 实际HAL main.c第110至133行选录 保留原行号的代码摘录}
\end{center}
\qsub{按照硬件使用顺序读代码}
第115、116行打开GPIOF和GPIOE的AHB1外设时钟。第119行先把PF9/\allowbreak{}PF10写为高电平，再切换为输出模式，使低电平有效的LED初始关闭；第122至126行设置输出模式、无上下拉和低输出速度。

第129至132行把PE4改为上拉输入。输入端悬空时由上拉保持高电平，按下按键接地后读到GPIO\_\allowbreak{}PIN\_\allowbreak{}RESET。GPIO速度控制的是输出边沿驱动特性，GPIO\_\allowbreak{}SPEED\_\allowbreak{}FREQ\_\allowbreak{}LOW不是软件翻转频率。

\begingroup\small
\begin{tabularx}{\linewidth}{|>{\hsize=0.419\hsize}Y|>{\hsize=1.186\hsize}Y|>{\hsize=1.395\hsize}Y|}\hline
\rowcolor{tablehead}
\textbf{目的} & \textbf{本工程使用的HAL函数} & \textbf{对应含义} \\ \hline
初始化 & HAL\_\allowbreak{}GPIO\_\allowbreak{}Init & 把结构体配置写入GPIO寄存器 \\ \hline
读引脚 & HAL\_\allowbreak{}GPIO\_\allowbreak{}ReadPin & 读取输入数据位 \\ \hline
写引脚 & HAL\_\allowbreak{}GPIO\_\allowbreak{}WritePin & 通过置位/\allowbreak{}复位机制输出指定电平 \\ \hline
翻转输出 & HAL\_\allowbreak{}GPIO\_\allowbreak{}TogglePin & 根据当前输出状态改变指定引脚 \\ \hline
\end{tabularx}
\endgroup\par

\clearpage
\pdfbookmark[0]{第8题 扩展练习四 HAL主循环与预期现象}{page15}
\qtitle{第8题 扩展练习四 HAL主循环与预期现象}
\begin{center}
\includegraphics[width=6.3in]{assets/图10_HAL主循环.png}\par
{\footnotesize 图10 实际HAL main.c第38至62行选录 省略注释与空行后保留原行号}
\end{center}
初始化次序为HAL\_\allowbreak{}Init → SystemClock\_\allowbreak{}Config → MX\_\allowbreak{}GPIO\_\allowbreak{}Init。进入循环后，PF9先翻转一次，调用HAL\_\allowbreak{}Delay(500)等待，再读取PE4：低电平时将PF10写低点亮LED1，高电平时写高关闭。

\qsub{由代码能够得出的结论}
\begingroup\small
\begin{tabularx}{\linewidth}{|>{\hsize=0.465\hsize}Y|>{\hsize=1.279\hsize}Y|>{\hsize=1.256\hsize}Y|}\hline
\rowcolor{tablehead}
\textbf{观察对象} & \textbf{按代码与配置推导的预期} & \textbf{不能直接推出的结果} \\ \hline
PF9 LED0 & 名义上约每500 ms翻转一次；一个完整亮灭周期约1 s & 不能据编译日志证明实际波形频率或LED已闪烁 \\ \hline
PE4与PF10 & 每轮延时后采样按键，按下时PF10输出低 & 不是中断即时响应；短按可能漏读，且没有消抖 \\ \hline
SysTick时基 & HAL采用1 ms中断时基；168 MHz时LOAD应为167999 & 未取得运行中的寄存器读数或示波器测量 \\ \hline
\end{tabularx}
\endgroup\par
HAL\_\allowbreak{}Delay依赖SysTick\_\allowbreak{}Handler调用HAL\_\allowbreak{}IncTick来推进uwTick。HAL\_\allowbreak{}RCC\_\allowbreak{}ClockConfig切频后会重建时基；不能再把普中SPL中会重装并关闭SysTick的delay\_\allowbreak{}ms混入同一HAL工程，否则可能破坏HAL超时和延时。

这些现象是源码推导；实际按键响应、LED频率和定时误差需在具备硬件后另行测量。


\clearpage
\pdfbookmark[0]{第8题 扩展练习四 SPL GPIO与时钟对照}{page16}
\qtitle{第8题 扩展练习四 SPL GPIO与时钟对照}
\begin{center}
\includegraphics[width=6.3in]{assets/图12_SPL_GPIO初始化.png}\par
{\footnotesize 图11 实际SPL SysTick工程led.c第9至22行 GPIO初始化源码}
\end{center}
\begin{center}
\includegraphics[width=6.3in]{assets/图11_SPL时钟主循环.png}\par
{\footnotesize 图12 实际SPL时钟工程main.c第57至68行 主循环源码}
\end{center}
SPL的RCC\_\allowbreak{}AHB1PeriphClockCmd、GPIO\_\allowbreak{}Init、GPIO\_\allowbreak{}ResetBits和GPIO\_\allowbreak{}SetBits分别对应时钟使能、配置、输出低和输出高。读取输入电平使用GPIO\_\allowbreak{}ReadInputDataBit；例如教师CH5EX02读取GPIOF的GPIO\_\allowbreak{}Pin\_\allowbreak{}6，对应HAL\_\allowbreak{}GPIO\_\allowbreak{}ReadPin。PF6是教师板按键，不能替换成普中板的PE4。两种库都要先开时钟、再设模式、再读写。

SPL时钟例程在main中调用RCC\_\allowbreak{}HSE\_\allowbreak{}Config(8,336,4,7)。HSE=8 MHz时，SYSCLK=8÷8×336÷4=84 MHz，HCLK=84 MHz，PCLK1=21 MHz，PCLK2=42 MHz。不能把此工程写成HAL工程的168/\allowbreak{}42/\allowbreak{}84 MHz。

delay(6000000)是空循环忙等，时长还受优化和指令执行影响，不是精确的500 ms。该例程重新配置时钟后未调用SystemCoreClockUpdate，软件时钟变量仍可能保留原168 MHz值；再次根据它配置时基前必须更新。


\clearpage
\pdfbookmark[0]{第8题 扩展练习四 SPL SysTick的使用}{page17}
\qtitle{第8题 扩展练习四 SPL SysTick的使用}
第二个SPL工程在SystemInit内配置168 MHz，在main中调用SysTick\_\allowbreak{}Init(168)。它采用HCLK/\allowbreak{}8作为SysTick计数时钟，即21 MHz，每微秒21个计数。此处的168是传入的MHz数值，必须与实际HCLK一致。

\begin{Verbatim}[fontsize=\scriptsize,breaklines=true]
/* Actual SPL SysTick User/main.c, lines 21-34 */
int main()
{
    SysTick_Init(168);
    LED_Init();
    while(1)
    {
        LED1=0;
        LED2=1;
        delay_ms(500);
        LED1=1;
        LED2=0;
        delay_ms(500);
    }
}
\end{Verbatim}
上段保留实际执行语句，仅调整缩进并省略源码中的延时精度注释。该例配套LED驱动初始化PF9与PF10；主程序交替改变两路LED状态，调用delay\_\allowbreak{}ms(500)后再进入另一种状态。

\qsub{沿着调用链检查延时}
delay\_\allowbreak{}ms把较长等待分成540 ms的小段，再调用delay\_\allowbreak{}nms。delay\_\allowbreak{}nms写入LOAD，清零VAL，打开计数器，轮询CTRL的COUNTFLAG，结束后关闭计数器并清零VAL。因此它是占用CPU等待的轮询式延时，不是HAL的全局毫秒计数。

\begingroup\small
\begin{tabularx}{\linewidth}{|>{\hsize=0.465\hsize}Y|>{\hsize=1.279\hsize}Y|>{\hsize=1.256\hsize}Y|}\hline
\rowcolor{tablehead}
\textbf{项目} & \textbf{SPL课程例程} & \textbf{HAL最小工程} \\ \hline
时钟来源 & SysTick=HCLK/\allowbreak{}8；168 MHz下为21 MHz & 默认SysTick=HCLK；168 MHz下为168 MHz \\ \hline
等待方法 & 读取COUNTFLAG直到计满 & SysTick中断累加uwTick，HAL\_\allowbreak{}Delay等待差值 \\ \hline
结束状态 & delay\_\allowbreak{}nms会关闭SysTick计数器 & 1 ms周期中断时基继续运行 \\ \hline
频率改变 & 须重新核对SysTick\_\allowbreak{}Init参数与延时系数 & 正确执行HAL切频流程后重建时基 \\ \hline
\end{tabularx}
\endgroup\par
源码把LOAD直接写为nms×fac\_\allowbreak{}ms，周期分析应考虑LOAD+1、启动和软件开销。课程注释中的“精确延时”不等于已测得的误差保证；本次只据源码说明机制，不把500 ms写成实测值。


\clearpage
\pdfbookmark[0]{第8题 CubeMX生成工程与时钟参数}{page18}
\qtitle{第8题 CubeMX生成工程与时钟参数}
clock168\_\allowbreak{}gpio.ioc选择STM32F407ZGT6，HSE为8 MHz，PF9为GPIO输出。生成环境是CubeMX 6.18.1和STM32CubeF4 V1.28.3，目标工具链为MDK-ARM V5。实际生成的SystemClock\_\allowbreak{}Config如下。

\begin{center}
\includegraphics[width=6.3in]{assets/图13_CubeMX时钟配置.png}\par
{\footnotesize 图13 CubeMX生成的Src/\allowbreak{}main.c第121至146行选录 实际代码摘录}
\end{center}
\begingroup\small
\begin{tabularx}{\linewidth}{|>{\hsize=0.651\hsize}Y|>{\hsize=1.023\hsize}Y|>{\hsize=1.326\hsize}Y|}\hline
\rowcolor{tablehead}
\textbf{参数} & \textbf{代码设置} & \textbf{对应频率或条件} \\ \hline
PLL M N P Q & 8、336、2、7 & VCO输入1 MHz，VCO输出336 MHz \\ \hline
SYSCLK /\allowbreak{} HCLK & PLLCLK、AHB不分频 & 168 MHz /\allowbreak{} 168 MHz \\ \hline
PCLK1 /\allowbreak{} PCLK2 & APB1四分频、APB2二分频 & 42 MHz /\allowbreak{} 84 MHz \\ \hline
Flash与电压等级 & FLASH\_\allowbreak{}LATENCY\_\allowbreak{}5、Scale 1 & 按相应供电条件支持168 MHz工作 \\ \hline
\end{tabularx}
\endgroup\par
实际main调用HAL\_\allowbreak{}Init、SystemClock\_\allowbreak{}Config和MX\_\allowbreak{}GPIO\_\allowbreak{}Init后进入空while(1)。gpio.c将PF9初始化为低电平，按参考板LED低有效接法预期为常亮；循环未实现动态翻转或按键处理。这与视频中的其他演示工程、HAL最小工程是不同版本。

时钟获取也要区分：HAL\_\allowbreak{}RCC\_\allowbreak{}GetSysClockFreq依据RCC配置计算SYSCLK；HAL\_\allowbreak{}RCC\_\allowbreak{}GetHCLKFreq返回SystemCoreClock缓存。直接改寄存器后若未维护缓存，后者不能自动代表当前实际时钟。


\clearpage
\pdfbookmark[0]{第8题 Keil历史编译记录}{page19}
\qtitle{第8题 Keil历史编译记录}
下表摘录保存的AC5原生.build\_\allowbreak{}log.htm。六个目标的记录均为0 Error(s)、0 Warning(s)，其中前四项属于原作业参考工程，后两项属于教师CH5EX01原工程。它们是既有构建证据，不是本轮修订源码的重编结果。大小单位为字节。

\begingroup\small
\begin{tabularx}{\linewidth}{|>{\hsize=1.442\hsize}Y|>{\hsize=0.679\hsize}Y|>{\hsize=0.791\hsize}Y|>{\hsize=0.791\hsize}Y|>{\hsize=0.791\hsize}Y|>{\hsize=1.507\hsize}Y|}\hline
\rowcolor{tablehead}
\textbf{目标} & \textbf{Code} & \textbf{RO-data} & \textbf{RW-data} & \textbf{ZI-data} & \textbf{结果} \\ \hline
SPL时钟 & 1524 & 408 & 0 & 1632 & 0错 0警告 \\ \hline
SPL SysTick & 1272 & 424 & 4 & 1636 & 0错 0警告 \\ \hline
HAL最小工程 & 3906 & 450 & 16 & 1632 & 0错 0警告 \\ \hline
CubeMX生成 & 3142 & 442 & 16 & 1632 & 0错 0警告 \\ \hline
教师CH5EX01 STD & 1528 & 424 & 40 & 1648 & 0错 0警告 \\ \hline
教师CH5EX01 HAL & 3390 & 442 & 36 & 1636 & 0错 0警告 \\ \hline
\end{tabularx}
\endgroup\par
\begin{center}
\includegraphics[width=6.3in]{assets/图14_Keil编译结果局部.png}\par
{\footnotesize 图14 实际Keil编译完成窗口局部 上传的真实截图 HAL目标为0错误0警告}
\end{center}
图14裁自“Keil\_\allowbreak{}HAL工程编译完成”截图的Build Output区域，原图随附，显示Code=3906、链接及HEX生成信息，与HAL原生日志一致。另一个名为“Rebuild结果”的截图只显示Rebuild开始阶段，不将它当作完成证据。

\qsub{编译成功说明什么}
记录说明对应版本的源码在所列工具链、宏和依赖条件下完成编译、链接并生成映像。编译器不会自动验证外部晶振频率、PLL电气边界、板级接线或LED闪烁频率，因此编译通过不消除下一节所列配置问题。

记录采用Arm Compiler 5.06 update 7 build 960。Code、RO-data、RW-data和ZI-data含义不同，不能把Code列单独当作全部Flash或RAM占用。原始日志与修改后的源码必须分别保存，不能把旧日志作为新补丁通过编译的凭据。

\qsub{复核所需资料}
工程资料中保留四份原始文本日志和四份原生HTML日志，源码选集用于逐行阅读；完整工程入口见参考资料[6]。重建时应使用匹配的AC5与F4 DFP环境，并保留工程中的启动、系统、驱动和头文件依赖。


\clearpage
\pdfbookmark[0]{第8题 教师CH5EX01工程对照}{page20}
\qtitle{第8题 教师CH5EX01工程对照}
CH5EX01是时钟讨论题指定的参考例。教师HAL与STD均以PE3、PE4为LED输出，PB0至PB7为流水灯输出；CH5EX02还使用PF6等按键。普中参考板则把PF9、PF10用于LED，PE4用于按键。两者的PE4方向相反，移植前必须按实际板卡核对。

\begingroup\small
\begin{tabularx}{\linewidth}{|>{\hsize=0.628\hsize}Y|>{\hsize=1.186\hsize}Y|>{\hsize=1.186\hsize}Y|}\hline
\rowcolor{tablehead}
\textbf{项目} & \textbf{CH5EX01 HAL原件} & \textbf{CH5EX01 STD原件} \\ \hline
HSE与PLL & 8 MHz；M8 N336 P2 Q4 & 25 MHz；M25 N336 P2 Q7 \\ \hline
主频与总线 & 168 /\allowbreak{} 168 /\allowbreak{} 42 /\allowbreak{} 84 MHz & 25 MHz实际输入时同左 \\ \hline
SysTick & 默认HCLK；1 ms中断 & HCLK/\allowbreak{}100；10 ms中断 \\ \hline
等待函数 & HAL\_\allowbreak{}Delay(1000) & Delay(10)，即约100 ms \\ \hline
PE3翻转 & 每轮翻转，约0.5 Hz完整周期 & 每5轮翻转，稳态约1 Hz \\ \hline
启动路径 & SystemInit后，main中配置PLL & SystemInit内调用SetSysClock \\ \hline
\end{tabularx}
\endgroup\par
\qsub{两个需要修正的时钟条件}
HAL原件的PLLQ=4使336 MHz VCO经Q支路输出84 MHz，超过DS8626 Rev12表36给出的75 MHz上限。保持CPU为168 MHz时，应在独立修订版把Q改为7，得到48 MHz，并同步IOC；不能把原配置整体表述为符合电气规格。

STD原件按25 MHz晶振设置HSE\_\allowbreak{}VALUE=25000000、PLL\_\allowbreak{}M=25。若直接换到8 MHz板，PLL输入只有0.32 MHz，低于规定范围。53.76 MHz只是代入公式的理想结果，实际不保证锁定或运行。8 MHz移植版须同时改HSE\_\allowbreak{}VALUE与PLL\_\allowbreak{}M，而不是只改一个软件常量。

教师STD的HSE失败分支没有恢复168 MHz配置，启动流程可能留在HSI 16 MHz；SystemCoreClock全局初值也可能过期。HAL配置失败会进入Error\_\allowbreak{}Handler，不能描述成无缝切换并继续原频率运行。


\clearpage
\pdfbookmark[0]{第8题 课件 视频与工程版本辨析}{page21}
\qtitle{第8题 课件 视频与工程版本辨析}
同一个教学主题出现了不同演示版本。表中频率按实际画面或源码计算，不能把视频、PPT和RAR中的参数统一成一个数值。这里的时间戳用于回到视频核对，不表示板上测得的时间。

\begingroup\small
\begin{tabularx}{\linewidth}{|>{\hsize=0.907\hsize}Y|>{\hsize=1.0\hsize}Y|>{\hsize=1.093\hsize}Y|}\hline
\rowcolor{tablehead}
\textbf{材料与定位} & \textbf{所见配置} & \textbf{对应结论} \\ \hline
CubeMX新建工程视频 10分20秒 & HSE25 /\allowbreak{} M20 /\allowbreak{} N268 /\allowbreak{} P2 & SYSCLK 167.5 MHz；演示接受接近168的配置 \\ \hline
HAL GPIO视频 4分15秒 & HSE8 /\allowbreak{} M8 /\allowbreak{} N336 /\allowbreak{} P2 & SYSCLK 168 MHz \\ \hline
HAL GPIO视频 18分15秒 & Delay100；每5轮翻转PE3 & 约500 ms翻转，完整周期约1 Hz \\ \hline
SPL GPIO视频 19分45秒及22分15秒 & 10 ms tick；Delay20；每轮翻转 & 约200 ms翻转，完整周期约2.5 Hz \\ \hline
第5章课件 例5-1 & PE3心跳1 Hz，PB依次循环 & 功能要求，不能代替版本实际实现 \\ \hline
教师HAL RAR & HAL\_\allowbreak{}Delay(1000)后循环翻转 & 约1000 ms翻转，完整周期约0.5 Hz \\ \hline
\end{tabularx}
\endgroup\par
\qsub{怎样使用这些材料}
先确定采用哪份源工程，再检查该版本的晶振、PLL参数、引脚和SysTick。若要完成1 Hz心跳加100 ms流水灯，可按100 ms推进流水，再以5步为一组翻转心跳，或使用独立的毫秒时间差调度。不能只把一个Delay改成500，就认为所有输出节拍同时满足。

GPIO视频中“低电平点亮”的接法与输出初始高电平相对应。输出模式和速度是硬件电气设置，Very High不代表LED软件翻转频率。视频里的工程建立、编译或下载配置画面也不能自动作为当前实体板已经运行的证据。

SPL新建工程视频8分10秒的结果为0错误、1警告；它与上传的教师工程0错0警历史日志属于不同记录，分别保留。


\clearpage
\pdfbookmark[0]{第8题 仿真记录与实验结论}{page22}
\qtitle{第8题 仿真记录与实验结论}
HAL工程已进入µVision的Simulation调试会话。下图显示程序在HAL\_\allowbreak{}Init启用Flash缓存附近停止，命令窗口报告对0x40023C00的读取访问受限。两幅图均裁自同一张实际IDE截图，未重绘界面或报错内容。

\begin{center}
\includegraphics[width=6.3in]{assets/图15_HAL初始化停止位置.png}\par
{\footnotesize 图15 实际仿真截图中的HAL\_\allowbreak{}Init停止位置 源码窗口局部}
\end{center}
\begin{center}
\includegraphics[width=5.5in]{assets/图16_仿真错误信息.png}\par
{\footnotesize 图16 同一仿真截图的命令窗口 error 65访问限制}
\end{center}
0x40023C00属于STM32F407的Flash接口寄存器区。本次会话在初始化阶段就受阻，因而没有形成PLL、GPIO或HAL时基持续运行的证据。这个结果是仿真会话的访问限制，不能解释为实体芯片或硬件电路失败。

SysTick属于Cortex-M内核部件，仅凭这条Flash寄存器报错，不能进一步断言所有内核功能都不支持。最稳妥的记录是：工程可构建，调试会话可启动，但本次外设运行验证未完成。

\begingroup\small
\begin{tabularx}{\linewidth}{|>{\hsize=1.008\hsize}Y|>{\hsize=0.992\hsize}Y|}\hline
\rowcolor{tablehead}
\textbf{已完成} & \textbf{本次未形成的证据} \\ \hline
四工程源码、配置与启动链阅读；实际Keil构建日志核对 & 实体板烧录、LED与按键实测、示波器或逻辑分析仪频率测量 \\ \hline
HAL实际IDE截图；仿真停止位置与报错记录 & 运行态时钟寄存器、SysTick持续计数及GPIO波形 \\ \hline
\end{tabularx}
\endgroup\par
原题第8题为自选扩展练习，未把烧录或上板设为提交前提。本次按现有资料完成书面分析与已有工程记录复核，保留验证边界，不补写不存在的实物现象。


\clearpage
\pdfbookmark[0]{参考资料与工程来源}{page23}
\qtitle{参考资料与工程来源}
下列资料用于核对本作业所涉及的章节，不代表通读全部手册。查阅日期为2026年10月4日。正文以方括号编号指向相应资料。

\href{https://www.st.com/resource/en/datasheet/stm32f407zg.pdf}{[1] ST DS8626 Rev12 STM32F405xx和STM32F407xx数据手册}\par
核对Features、LQFP144引脚图、表7引脚表、存储器映射、时钟电气条件及NRST推荐电路。随附PDF：sources/\allowbreak{}STM32F407ZG\_\allowbreak{}DS8626.pdf。

\href{https://www.st.com/resource/en/reference_manual/dm00031020-stm32f405415-stm32f407417-stm32f427437-and-stm32f429439-advanced-armbased-32bit-mcus-stmicroelectronics.pdf}{[2] ST RM0090 STM32F4参考手册}\par
核对F405/\allowbreak{}407适用的系统结构、存储器映射、RCC和GPIO章节。RM0090同时覆盖多种F4器件，必须先看章节适用范围，再使用总线与PLL参数。

\href{https://www.st.com/resource/en/programming_manual/pm0214-stm32-cortexm4-mcus-and-mpus-programming-manual-stmicroelectronics.pdf}{[3] ST PM0214 Rev10 Cortex-M4编程手册}\par
核对寄存器组、存储器映射、异常模型、向量表、NVIC、SCS和SysTick章节。相应PDF保存在sources目录。

\href{https://www.st.com/resource/en/datasheet/stm32l496ve.pdf}{[4] ST DS11585 Rev20 STM32L496xx数据手册}\par
核对Features、具体器件对照表及订货型号说明。随附PDF：sources/\allowbreak{}STM32L496VE\_\allowbreak{}DS11585.pdf。

\href{https://www.st.com/en/microcontrollers-microprocessors/stm32l496ve.html}{[5] ST STM32L496VE器件页面}\par
用于交叉核对STM32L496VET6的512 KB Flash、LQFP100和温度等级。


\clearpage
\pdfbookmark[0]{参考资料与版本定位}{page24}
\qtitle{参考资料与版本定位}
\href{https://drive.google.com/drive/folders/1PP6caucxd0rbkspc_7IDALtIaz1bU5kz}{[6] 四份实际工程及本机编译记录}\par
本次代码摘录、工程分组、构建表和实际界面图均据此目录中的对应文件。图8至13是实际源码排版摘录，不是IDE截图；图14至16是实际IDE截图的局部。

\href{https://drive.google.com/file/d/16aZsBIdPqjkYs0ldyzE22hmRR2uYA8LM/view}{[7] 2026年10月4日本机工程与工具验证记录}\par
环境和过程以该记录的后续补充段及原始工程文件交叉核对。较早的“CubeF4未安装”状态已被后续生成与构建成功记录更新。芯片参数、引脚和函数行为以对应手册、原理图和源码为准。

随附Word正文和表格可编辑；LaTeX源文件与配图可重新编译为PDF。工程资料中的源码为阅读选集，完整可构建工程请使用[6]中的原始目录。

\qsub{补充课程来源}
教师PPT：第3章Cortex-M4内核体系结构；第4章STM32芯片应用基础；第5章GPIO与SysTick。总线矩阵核对RM0090 Rev19第2.1节第59页；PLL限制核对DS8626 Rev12表36；SysTick重载核对Arm DUI0553B第4.4.2节。

\href{https://drive.google.com/file/d/1SiaXWNtQBBjrw8Dl70PK68exQ3V929ve/view}{教师CH5EX01 HAL原工程}\par
\href{https://drive.google.com/file/d/1BVRke_0FVb91vsfTokTQznbfqVBA4-GZ/view}{教师CH5EX01 STD原工程}\par
视频定位采用课程提供的SPL新建工程、CubeMX新建工程、HAL GPIO及SPL GPIO演示文件。相关参数须与同版本源码对应；不能用课程画面冒充本轮软件运行截图。

\end{document}
